\section{Contexte et objectif}

\subsection{Position dans le depot}

Le dossier \path{full_soft/} contenait une pile logicielle autonome plus ancienne que la pile APEX. Ce dossier n'est plus present dans \path{main}: il est considere comme une reference externe historique. La pile analysee ne reposait pas sur ROS 2; elle utilisait des modules Python pour lire les capteurs, calculer une direction et une vitesse, commander les actionneurs et afficher des visualisations.

Cette architecture semble avoir ete pensee pour une exploitation directe sur Raspberry Pi, avec une logique de developpement pedagogique: installation, connexion au Raspberry Pi, configuration reseau, VNC, VSCode, scripts de dependances et documentation hardware.

\subsection{Lien avec la Voiture Blue}

La documentation materielle historique mentionnait plusieurs voitures et composants, dont la Voiture Blue. Cette reference peut donc servir a comprendre l'ancien controle materiel: LiDAR, camera, moteur, direction, ultrason, batterie et vitesse.

En revanche, d'apres l'etat actuel du depot, l'usage recommande pour la Voiture Blue est APEX/ROS 2. L'utilisation de cette pile Python avec la Voiture Blue doit donc etre consideree comme une approche historique, de transition ou de reference externe, plutot que comme le chemin principal actuel.

\subsection{Objectif du rapport}

Ce rapport vise a:

\begin{itemize}
  \item decrire l'organisation Python de l'ancien miroir \path{full_soft/};
  \item expliquer la logique d'execution et les interfaces;
  \item faire un bilan d'utilisation pour la Voiture Blue;
  \item identifier les problemes confirmes et les risques potentiels;
  \item distinguer clairement cette pile de la pile ROS 2 APEX.
\end{itemize}
